% Clase 17 --- La ley de Ampère: qué corrientes cuentan y con qué signo (§3.1-3.2).
% "Si tomamos la corriente orientada para allá, las corrientes que pasen de abajo
%  para arriba las vamos a tomar positivas, y las que pasen de arriba para abajo,
%  negativas. Hay que orientar las corrientes de acuerdo con la regla de la mano
%  derecha en función de cómo orientamos la curva."
% (a) el enunciado y su convención; (b) el caso donde la ley sirve de verdad,
% que es cuando la simetría permite sacar B de la integral.
\begin{center}
\begin{tikzpicture}[scale=1]

  % =============== (a) la curva y las corrientes ===============
  \begin{scope}
    % curva cerrada irregular, orientada
    \draw[figline, line width=1.1pt]
      (-1.9,0.1) .. controls (-1.7,1.6) and (-0.3,1.9) .. (0.6,1.5)
      .. controls (1.6,1.05) and (2.0,0.2) .. (1.6,-0.8)
      .. controls (1.2,-1.7) and (-0.6,-1.9) .. (-1.5,-1.2)
      .. controls (-2.0,-0.8) and (-2.05,-0.5) .. (-1.9,0.1);
    \draw[figvec, line width=1.1pt] (-1.83,0.65) -- (-1.72,1.15);
    \node[figlbl, figblue, anchor=east] at (-1.9,0.95) {$C$};
    \node[figlblaux, anchor=east, align=right] at (-1.95,1.55)
      {$C$ orientada; el pulgar\\[-0.25em] da el signo de $I$};

    % corrientes que la atraviesan
    \node[figred, scale=1.15] at (-0.55,0.55) {$\odot$};
    \node[figlbl, figred, anchor=south] at (-0.55,0.72) {$+I_1$};
    \node[figred, scale=1.15] at (0.6,0.05) {$\otimes$};
    \node[figlbl, figred, anchor=west] at (0.75,0.05) {$-I_2$};
    \node[figred, scale=1.15] at (-0.7,-0.95) {$\odot$};
    \node[figlbl, figred, anchor=north] at (-0.7,-1.12) {$+I_3$};

    % una corriente exterior, que no cuenta
    \node[figgray, scale=1.15] at (2.5,1.35) {$\odot$};
    \node[figlblaux, anchor=west, align=left] at (2.65,1.3)
      {$I_4$ no atraviesa $C$:\\[-0.25em] no entra en\\[-0.25em] la suma};

    \node[figlbl, anchor=north, align=center] at (0.3,-2.35)
      {(a) $\displaystyle\oint_C \vec B\cdot d\vec l =
       \mu_0 (I_1 - I_2 + I_3)$};
  \end{scope}

  % =============== (b) el hilo infinito ===============
  \begin{scope}[xshift=7.45cm]
    % la amperiana circular
    \draw[figline, line width=1.1pt] (0,0) circle (1.55);
    \draw[figvec, line width=1.1pt] (100:1.55) arc (100:80:1.55);
    \node[figlbl, figblue, anchor=south] at (0,1.6) {$C$};

    % el hilo, de punta
    \node[figred, scale=1.25] at (0,0) {$\odot$};
    \node[figlbl, figred, anchor=north east] at (-0.12,-0.1) {$I$};

    % B tangente y de módulo constante
    \foreach \ang in {0,45,90,135,180,225,270,315} {
      \draw[figvecres, line width=0.9pt]
        (\ang:1.55) -- ($(\ang:1.55)+({\ang+90}:0.6)$);
    }
    \node[figlbl, figred, anchor=west] at (1.15,-1.35) {$\vec B$};

    % el radio
    \draw[figvecaux] (0,0) -- (-1.55,0);
    \node[figlblaux, anchor=south] at (-0.8,0.04) {$R$};

    \node[figlbl, anchor=north, align=center] at (0,-2.35)
      {(b) $B\,(2\pi R) = \mu_0 I \Rightarrow
       B = \dfrac{\mu_0 I}{2\pi R}$};
  \end{scope}

\end{tikzpicture}\\[0.5em]
{\small\itshape La ley de Ampère juega para el magnetismo el papel que la de
Gauss para la electricidad, con la misma lógica y las mismas limitaciones. Del
lado izquierdo hay una circulación sobre una curva \textbf{cerrada}; del
derecho, sólo las corrientes que \textbf{atraviesan} esa curva. Las de afuera no
figuran, y eso no significa que no generen campo: lo generan, y de hecho
modifican $\vec B$ punto por punto sobre $C$; lo que ocurre es que su
contribución neta a la circulación se cancela. Los signos hay que tomarlos en
serio: una vez elegido el sentido de recorrido de $C$, el pulgar de la mano
derecha dice cuáles corrientes son positivas, y las que van al revés restan.
Como con Gauss, la ley vale siempre pero \emph{sirve} sólo cuando hay tanta
simetría que $B$ es constante sobre la curva y sale de la integral: en el hilo
infinito, las simetrías de rotación y de traslación obligan a que las líneas
sean círculos con $|\vec B|$ constante, y todo el cálculo de Biot--Savart de la
clase pasada se reduce a un renglón. Para un segmento finito no hay tal simetría
y hay que volver a Biot--Savart.}
\end{center}
